\documentclass[11pt,a4paper]{article}
\usepackage{turkos}

\begin{document}
\phaseheader{11}{Process API: fork, exec, wait and exit}{Implement the UNIX process model, and use the page-fault handler to make fork cheap and memory grow on demand}{70}{77}

\ataglance{3 weeks}{5}{Phase 10: user processes in their own address spaces, system calls, ELF loading, \code{copy_from_user}}{\code{fork}
with copy-on-write, \code{execve} with \code{argc}/\code{argv}, \code{waitpid}, \code{exit} with zombies and orphan
reparenting, \code{getppid}, \code{brk} with lazy allocation, a stack that grows on demand, minimal signals
(\code{kill}, \code{SIGSEGV}, \code{SIGINT}), and a user-space \code{init} process that starts and reaps other programs.}

\section{Why this phase}

UNIX creates processes in an unusual way, and it turned out to be one of its best ideas. Instead of one call that
starts a program, there are two. \code{fork} makes an exact copy of the calling process. \code{exec} replaces the
calling process's program with a new one. The parent then uses \code{wait} to learn when its child has finished.
\OSTEP chapter 5 explains why the split is so powerful: between \code{fork} and \code{exec}, the child can rearrange
its own environment (redirect its output to a file, connect a pipe, change directory) without the program it is about
to run knowing anything about it. Your shell in Phase 14 depends on exactly that.

A naive \code{fork} copies every page of the parent, only for \code{exec} to throw the copy away a moment later. The
fix is \textbf{copy-on-write} (COW): share the pages and copy each one only when someone writes to it. COW is your
first use of a page fault as a \emph{feature} rather than an error, and the same idea gives you lazily allocated heaps
and stacks that grow on demand. After this phase, the page-fault handler is one of the most important functions in
the kernel.

\section{Concepts you need}

\subsection{The four calls and their contracts}

\textbf{\code{fork()}} creates a child that is a copy of the parent: same memory contents, same registers, same open
files later. It returns twice: the child's PID in the parent and 0 in the child. \textbf{\code{execve(path, argv,
envp)}} loads a new program into the current process, keeping its PID and open files; on success it never returns.
\textbf{\code{exit(status)}} ends the process, but it can't vanish entirely: its exit status must be kept until the
parent collects it, so it becomes a \textbf{zombie}. \textbf{\code{waitpid(pid, &status, options)}} blocks until a
child exits, collects its status and frees what is left of it. If a parent exits first, its children become
\textbf{orphans} and are adopted by \code{init} (PID 1), whose job includes reaping them (Figure~\ref{fig:tree}).

\begin{figure}[htbp]
\centering
\begin{tikzpicture}[proc/.style={draw=tkNavy!70, fill=tkLight, rounded corners=2pt, font=\sffamily\scriptsize, align=center, minimum width=2.2cm, minimum height=0.8cm},
  zomb/.style={proc, fill=tkGray!18, dashed}, level distance=1.35cm, sibling distance=3.0cm, edge from parent/.style={draw=tkNavy!70, -{Stealth[length=2mm]}}]
  \node[proc, fill=tkRed!10] {\textbf{init} (pid 1)\\reaps orphans}
    child {node[proc] {\textbf{tsh} (pid 5)\\the shell}
      child {node[proc] {\textbf{ls} (pid 9)\\\code{fork} + \code{exec}}}
      child {node[zomb] {\textbf{cat} (pid 8)\\zombie until \code{wait}}}}
    child {node[proc] {\textbf{daemon} (pid 7)\\parent exited:\\adopted by init}};
\end{tikzpicture}
\caption{A Turk-OS process tree in Phase 14. Every process except init has a parent that must eventually \code{wait} for it.}
\label{fig:tree}
\end{figure}

\subsection{Copy-on-write}

At \code{fork}, instead of copying the parent's pages, the kernel makes the child's page tables point to the same
frames and marks every writable page \emph{read-only} in \emph{both} processes, with a spare PTE bit (one of the AVL
bits 9--11) recording that it is really a COW page. Reading works as normal. The first write by either process causes
a page fault; the handler sees the COW bit, copies that one page to a fresh frame, maps the copy writable for the
faulting process, and returns. The CPU retries the write, which now succeeds (Figure~\ref{fig:cow}). A per-frame
\textbf{reference count} tells the handler how many address spaces share a frame: if it is already down to one, the
handler just makes the page writable again. \OSC\,\S10.3 describes COW, and \OSC\,\S10.2 describes the demand paging
it builds on.

\begin{figure}[htbp]
\centering
\begin{tikzpicture}[pt/.style={draw=tkNavy!70, fill=tkLight, minimum width=2.4cm, minimum height=0.55cm, font=\sffamily\scriptsize},
  fr/.style={draw=tkNavy!70, minimum width=1.9cm, minimum height=0.9cm, font=\sffamily\scriptsize, align=center}]
  \node[font=\sffamily\small\bfseries, text=tkNavy] at (2.4,1.6) {right after \code{fork}};
  \node[pt] (a1) at (0,0.6) {parent PTE: R, COW};
  \node[pt] (b1) at (0,-0.6) {child PTE: R, COW};
  \node[fr, fill=tkAmber!15] (f1) at (4.3,0) {frame X\\refcount 2};
  \draw[tkarrow] (a1.east) -- (f1); \draw[tkarrow] (b1.east) -- (f1);
  \node[font=\sffamily\small\bfseries, text=tkNavy] at (11.0,1.6) {after the child writes};
  \node[pt] (a2) at (8.6,0.6) {parent PTE: R, COW};
  \node[pt] (b2) at (8.6,-0.6) {child PTE: R/W};
  \node[fr, fill=tkAmber!15] (f2) at (12.9,0.6) {frame X\\refcount 1};
  \node[fr, fill=tkGreen!12] (f3) at (12.9,-0.6) {frame Y (copy)\\refcount 1};
  \draw[tkarrow] (a2.east) -- (f2); \draw[tkarrow] (b2.east) -- (f3);
  \draw[tkarrow, tkRed, line width=0.9pt] (5.5,0) -- node[above, font=\sffamily\scriptsize, text=tkRed] {write fault} (7.1,0);
\end{tikzpicture}
\caption{Copy-on-write for one page. The parent's next write finds refcount 1 and only needs its PTE made writable.}
\label{fig:cow}
\end{figure}

\subsection{Faults as a feature}

\MOS\,\S3.6.2 lists the steps of page-fault handling in a real kernel, and \OSTEP chapters 21--22 go further: when
memory runs out, pages can be written to disk and brought back by a fault (swapping), which needs a replacement
policy such as the clock algorithm. Turk-OS uses faults for three things in this phase. \textbf{COW}: a write to a
present page marked COW. \textbf{Lazy heap}: \code{brk} only moves the heap's end, and a fault inside the heap maps a
zeroed page. \textbf{Stack growth}: a fault just below the stack maps a new page, up to a limit. Anything else from
user mode is a real error and ends the process with \code{SIGSEGV}. Swapping is a stretch goal once you have a disk.

\section{Reading plan}

\begin{readingplan}
\must & \OSTEP & Ch.~5 \emph{Interlude: Process API} (\code{fork}, \code{wait}, \code{exec}, why the split, redirection). Do its homework on Linux first & 37--48 & The exact semantics you must reproduce. \\
\must & \OSDI & \S4.7 \emph{Overview of the MINIX 3 Process Manager}: memory layout, \code{fork}/\code{exit}/\code{wait} (\S4.7.4), \code{exec} (\S4.7.5), \code{brk} (\S4.7.6), signals (\S4.7.7) & 420--447 & A complete UNIX process model, explained step by step. \\
\must & \OSDI & \S4.8.3--4.8.6 \emph{Implementation of fork, exit, wait, exec, brk, signal handling} & 455--471 & The real code for the same calls. \\
\must & \OSC & \S10.2 \emph{Demand Paging}; \S10.3 \emph{Copy-on-Write} & 392--401 & Page faults as the mechanism behind COW and lazy allocation. \\
\should & \OSC & \S3.3 \emph{Operations on Processes} & 116--123 & Creation, termination, zombies and orphans. \\
\should & \MOS & \S3.6 \emph{Implementation Issues} (page fault handling, instruction backup, locking pages, backing store) & 233--240 & What a page-fault handler really has to deal with. \\
\should & \OSTEP & Ch.~21--22 \emph{Beyond Physical Memory: Mechanisms}, \emph{Policies} & 227--254 & Swapping and replacement policies, for the stretch goal. \\
\could & \MOS & \S3.4 \emph{Page Replacement Algorithms} & 209--222 & Optimal, NRU, FIFO, second chance, clock, LRU, working set. \\
\could & \MOS & \S10.3 \emph{Processes in Linux} & 733--753 & \code{fork}, \code{clone} and \code{exec} in a production kernel. \\
\end{readingplan}

\begin{minixbox}
The MINIX process manager is a user-space server, and its handlers read almost like pseudo-code:
\code{do_fork}, \code{do_pm_exit} and \code{do_waitpid} in \texttt{servers/pm/forkexit.c (18400)}, \code{do_exec} in
\texttt{servers/pm/exec.c (18700)}, \code{do_brk} in \texttt{servers/pm/break.c (19300)}, and signal handling in
\texttt{servers/pm/signal.c (19500)}. Notice how \code{exec} builds the new stack with the argument strings and then
patches the pointers in it: you will do the same in Task~11.6. One difference is fundamental: MINIX 3 has no paging,
so its \code{fork} must copy the child's whole data segment immediately. With paging you can share and copy lazily.
\end{minixbox}

\section{Build it}

\task{Processes, parents and the trap frame}
Extend \code{struct task} with \code{ppid}, a list of children, \code{exit_code}, a wait queue for ``one of my children
exited'', \code{heap_start} and \code{brk}, and \code{tf}: a pointer to the \code{struct regs} that the current system
call entered with. Set \code{current->tf = r} at the top of the syscall handler; \code{fork}, \code{exec} and signals
all need to read or change the saved user registers. Put a limit on the number of processes (say 64), so that a fork
bomb fails with \code{-EAGAIN} instead of exhausting memory.

\task{Frame reference counts}
Add a reference count per physical frame to the PMM: a \code{uint16_t} array covers the 512\,MiB direct map with
256\,KiB. \code{pmm_ref(pa)} increments; \code{pmm_unref(pa)} decrements and frees the frame when the count reaches
zero. From now on, freeing a user page means unreferencing it.

\task{fork, the simple version}
Get \code{fork} right before making it fast. Create the child task and a new address space, and for every present
user page in the parent, allocate a frame, copy the page through the direct map, and map it at the same address with
the same flags. The difficult part is returning to user mode in the child. Copy the parent's trap frame to the top of
the child's kernel stack and set its EAX to 0. Below it, build a \code{switch_context} frame whose return address is
the exit path of \code{isr_common}, the instructions after \code{call isr_handler}. Give that point a label,
\code{isr_return}. When the child is first scheduled, it ``returns'' through that path and \code{iret}s into user mode,
at the instruction after \code{int 0x80}, with EAX = 0.

\begin{lst}{c}{kernel/proc/fork.c (the tricky part)}
extern void isr_return(void);               /* label in isr.s, just after "add esp, 4" */

static void setup_child_stack(struct task *child, const struct regs *parent_tf)
{
    char *top = (char *)child->kstack + KSTACK_SIZE;
    struct regs *tf = (struct regs *)(top - sizeof(struct regs));
    *tf = *parent_tf;                        /* same user registers as the parent ...  */
    tf->eax = 0;                             /* ... except fork() returns 0 in the child */

    uint32_t *sp = (uint32_t *)tf;
    *--sp = (uint32_t)isr_return;            /* switch_context's ret goes here */
    *--sp = 0;                               /* ebp */
    *--sp = 0;                               /* ebx */
    *--sp = 0;                               /* esi */
    *--sp = 0;                               /* edi */
    child->esp = (uint32_t)sp;
    child->tf  = tf;
}
\end{lst}
\verify{A test where the parent sets a variable to 1 and the child sets it to 2 shows each process keeping its own
value, and \code{fork} returns the child's PID in the parent.}

\task{Copy-on-write}
Replace the copying loop. For every present user PTE in the parent: if it is writable, clear R/W and set your
\code{PG_COW} bit (\code{0x200}) in the parent's PTE; copy the PTE into the child; call \code{pmm_ref} on the frame.
When the loop is done, reload CR3 in the parent: you changed many of its live mappings, and its TLB still holds the
old writable ones. Then extend the page-fault handler.

\begin{lst}{c}{kernel/mm/fault.c (COW case)}
/* A write to a present page (error bits P and W set) below the kernel. The fault can come
 * from user mode or from the kernel writing through copy_to_user (CR0.WP is set). */
if ((r->err_code & 3) == 3 && addr < KERNEL_VMA) {
    pte_t *pte = vmm_get_pte(current->pgdir, addr);
    if (pte && (*pte & PG_COW)) {
        uintptr_t old = *pte & ~0xFFFu;
        if (pmm_refcount(old) == 1) {                    /* last user: just take it back */
            *pte = (*pte | PG_WRITE) & ~PG_COW;
        } else {
            uintptr_t copy = pmm_alloc_frame();
            if (!copy) { kill_current(SIGKILL); return; }
            memcpy(P2V(copy), P2V(old), PAGE_SIZE);
            *pte = copy | (*pte & 0xFFF & ~PG_COW) | PG_WRITE;
            pmm_unref(old);
        }
        __asm__ volatile ("invlpg (%0)" : : "r"(addr) : "memory");
        return;                                          /* the CPU retries the write */
    }
}
\end{lst}
\verify{Right after forking a process with a 1\,MiB array, the PMM free count drops only by a few frames (page tables
and the kernel stack). Each write to a new page of the array in the child costs exactly one more frame.}

\task{exit and waitpid}
\code{sys_exit} frees the user pages and page tables (unreferencing frames), gives every child to \code{init}, marks
the task as a zombie with its exit code, wakes the parent's wait queue, and schedules away. Keep the page directory
itself and the kernel stack for the parent to free: the exiting task is still running on both. \code{sys_waitpid}
looks for a zombie child that matches (\code{pid == -1} means any): if it finds one, it copies the status out with
\code{copy_to_user}, frees the child completely and returns its PID. If there are no children at all it returns
\code{-ECHILD}; with \code{WNOHANG} it returns 0; otherwise it sleeps on its wait queue and looks again.
\verify{1000 iterations of fork, exit and wait bring the PMM free count and the heap statistics back to exactly their
starting values. A child whose parent exits first is reaped by \code{init}.}

\task{execve with arguments}
The arguments live in the old address space, which you are about to throw away, so copy the path and the
\code{argv}/\code{envp} strings into kernel buffers first (with size limits). Build a complete new address space and
load the ELF into it. Only when everything has succeeded, switch to it, free the old one, and edit the trap frame:
EIP is the new entry point, the user ESP is the new stack pointer. That way a failed \code{exec} returns an error to a
process that still works. Lay out the new stack the way the System V i386 ABI does, and update \code{crt0}.

\begin{lst}{plain}{initial user stack built by execve (higher addresses at the top)}
0xC0000000  +-----------------------------------+
            | argument and environment strings  |  "hello\0" "-v\0" "PATH=/bin\0"
            +-----------------------------------+
            | NULL                              |
            | envp[0]                           |  pointers into the strings above
            | NULL                              |
            | argv[1], argv[0]                  |
ESP  -->    | argc                              |  _start reads this first
            +-----------------------------------+
\end{lst}

\begin{lst}{asm}{user/crt0.s (new version)}
_start:
    mov  eax, [esp]             ; argc
    lea  ebx, [esp + 4]         ; argv
    lea  ecx, [ebx + eax*4 + 4] ; envp = argv + argc + 1
    push ecx
    push ebx
    push eax
    call main                   ; main(argc, argv, envp)
    push eax
    call exit
\end{lst}
Until Phase 13 brings a file system, \code{execve} finds programs by name in the table of boot modules, so
\code{execve("/bin/hello", ...)} looks up the module called \code{hello}.
\verify{A program that prints its arguments shows exactly what its parent passed to \code{execve}. A failing
\code{execve} (unknown program) returns \code{-ENOENT} and the caller continues normally.}

\task{brk and a growing stack}
\code{sys_brk(addr)} returns the current end of the heap when \code{addr} is 0, and otherwise moves the end, refusing
to grow into the stack region; when the heap shrinks, unmap and unreference the pages beyond the new end. Don't map
anything when it grows: in the page-fault handler, a not-present fault inside \code{[heap_start, brk)} maps a zeroed
page. Similarly, a fault less than 8\,MiB below \code{0xC0000000} and near the current stack maps a new stack page.
\verify{A user program that \code{brk}s 64\,MiB but touches only three pages costs three frames.}

\task{Minimal signals}
Add a pending-signal bitmask to each task, and \code{sys_kill(pid, sig)}. Before returning to user mode (in
\code{isr_handler}, when the saved CS has RPL 3), check pending signals and apply the default action; for now every
signal Turk-OS supports terminates the process, with exit status $128 + \text{signal}$ so the parent can tell. Turn
user-mode exceptions into signals (\code{SIGSEGV} for page faults and \#GP, \code{SIGFPE} for divide errors,
\code{SIGILL} for invalid opcodes). Make Ctrl+C on the console send \code{SIGINT} to a designated foreground process;
Phase 14 formalises this with the terminal driver.

\task{init, and the tests}
Write \code{user/init.c}: the kernel starts it as PID 1 once the boot is done. For now it forks and execs the test
programs one after another and loops on \code{waitpid(-1, ...)} forever, which also reaps orphans. Add these tests to
\code{make test}.

\begin{center}
\small\renewcommand{\arraystretch}{1.25}
\begin{tabularx}{\linewidth}{@{}l X@{}}
\toprule
\tkth{Test} & \tkth{What must hold} \\
\midrule
\code{fork_values} & parent and child see different return values and keep separate copies of variables \\
\code{cow_frames} & forking a large process costs a few frames; each written page costs exactly one more \\
\code{fork_exec_wait} & a child execs a program that exits with 42; \code{waitpid} reports 42 \\
\code{orphans} & grandchildren of an exited process are reaped by init; no zombies remain \\
\code{no_leaks} & 1000 fork/exec/exit/wait cycles leave PMM and heap statistics unchanged \\
\code{fork_bomb} & forking without limit ends with \code{-EAGAIN}, and the system recovers when the processes exit \\
\code{signals} & \code{kill(child, SIGKILL)} ends a looping child; \code{waitpid} reports status $128+9$ \\
\bottomrule
\end{tabularx}
\end{center}

\section{Testing and debugging}

\begin{pitfalls}
The child crashes as soon as it runs & Its kernel stack doesn't match what \code{isr_return} pops, or the trap frame wasn't copied & Check in GDB that ESP at \code{isr_return} equals the address of the child's \code{struct regs}. \\
Parent and child see each other's writes after \code{fork} & The parent's TLB still holds writable entries for pages now marked COW & Reload CR3 in the parent after the COW loop. \\
A write by the kernel to user memory corrupts another process & CR0.WP not set, so the kernel wrote straight through a COW page & Set CR0.WP (Phase 6); the fault handler then copies the page for the kernel too. \\
Free memory shrinks a little with every fork & A reference count isn't dropped somewhere: exit, exec's old address space, or a COW copy & Run \code{no_leaks}; print the counts before and after one cycle. \\
Zombies pile up & Nobody waits: the shell or init doesn't reap & init loops on \code{waitpid(-1)}; later the shell reaps background jobs. \\
A failed \code{exec} leaves a broken process & The old address space was freed before the load succeeded & Build first, swap last. \\
\end{pitfalls}

\section{Milestone}

\begin{milestonebox}[The UNIX process model]
\begin{checklist}
  \item \code{fork} with copy-on-write, \code{execve} with arguments, \code{waitpid}, \code{exit}, \code{getpid}, \code{getppid} and \code{brk} work.
  \item The page-fault handler resolves COW, lazy heap and stack growth, and turns everything else into \code{SIGSEGV}.
  \item \code{init} runs as PID 1, starts programs and reaps every orphan.
  \item All seven tests pass, including 1000 cycles without leaks.
\end{checklist}
\end{milestonebox}

\begin{questionbox}
\begin{enumerate}
  \item Why does UNIX separate \code{fork} and \code{exec}? Give an example the split makes easy.
  \item How does one call to \code{fork} return two different values? Point to the line in your kernel that does it.
  \item Walk through copy-on-write for one page, including every TLB flush and reference count change.
  \item What is a zombie, what is an orphan, and who cleans each up?
  \item Why must \code{execve} copy its arguments into the kernel before building the new address space?
  \item What does \code{exec} keep from the old process, and what does it replace?
  \item Why does a lazy \code{brk} make programs that over-allocate cheap? What is the risk?
\end{enumerate}
\end{questionbox}

\section{Going further}

Once Phase 12 gives you a disk, add swapping: when the PMM runs low, pick a victim page with the clock algorithm
(\OSTEP chapter 22, \MOS\,\S3.4.5), write it to a swap partition, mark the PTE not present with the disk slot in the
free bits, and read it back on the next fault. Add user-level signal handlers with \code{sigaction} and a
\code{sigreturn} trampoline on the user stack. Or implement threads inside a process with a \code{clone}-style call
that shares the address space.

\nextphase{12}{Storage: Block Devices, ATA Driver and Buffer Cache}
\booklegend
\end{document}
